\begin{lemma}
\label{editorial-source-lemma-geometrically-connected}
Let $k$ be a field.
Let $R$ be a $k$-algebra.
The following are equivalent
\begin{enumerate}
\item for every field extension $k'/k$ the
spectrum of $R \otimes_k k'$ is connected, and
\item for every finite separable field extension $k'/k$ the
spectrum of $R \otimes_k k'$ is connected.
\end{enumerate}
\end{lemma}

\begin{proof}
Condition (1) implies (2).
Assume (2); its identity field
extension says that $\Spec(R)$
is connected and in particular
nonempty, so $R\neq0$.
Every scalar extension considered
below is nonzero by extension
of a vector-space basis.
For these nonzero rings, Lemma
\ref{algebra-lemma-characterize-spec-connected}
identifies connectedness with
absence of nontrivial idempotents.
The nonzero condition cannot be
omitted from that criterion:
the zero ring has no nontrivial
idempotents but has empty,
nonconnected spectrum under
Topology, Definition
\ref{topology-definition-connected-components}.

Let $k_s/k$ be the original
separable algebraic closure.
Every finite list of its
elements lies in a finite
separable subextension $E/k$.
The original coefficient maps
give the injective directed union
$$
R\otimes_k k_s=\bigcup_E R\otimes_k E.
$$
An idempotent in this union
lies in one such tensor, and
its equation is already true
there by injectivity. By (2)
it equals $0$ or $1$ there,
hence in the union. This
proves connectedness over $k_s$.

For any original field $L/k$,
choose a prime $\mathfrak a$
of the nonzero ring $L\otimes_k k_s$
and put
$$
F=\operatorname{Frac}((L\otimes_k k_s)/\mathfrak a).
$$
Its original field embeddings
$a\mapsto[a\otimes1]$ and
$b\mapsto[1\otimes b]$ are
injective and agree on $k$.
The scalar comparison maps
$$
(R\otimes_k k_s)\otimes_{k_s}F\longrightarrow R\otimes_k F,
\qquad(r\otimes a)\otimes b\longmapsto r\otimes ab,
$$
and $r\otimes b\mapsto(r\otimes1)\otimes b$
are inverse by balancing.
The source has connected
spectrum by Lemma
\ref{editorial-source-lemma-separably-closed-connected},
since both $R\otimes_k k_s$
and the field $F$ do.
The injection
$R\otimes_k L\to(R\otimes_k L)\otimes_L F\simeq R\otimes_k F$
reflects nontrivial idempotents,
again by a vector-space basis
and the displayed scalar maps.
Thus $R\otimes_k L$ is nonzero
and has no nontrivial idempotents.
It is connected, proving (1).

The same argument proves that
connectedness after just the
separable closure is an equivalent
test. One may instead use the
algebraic closure $\overline k$:
the extension $\overline k/k_s$
is purely inseparable, or the
identity in characteristic zero.
The original injection
$R\otimes_k k_s\to R\otimes_k\overline k$
is a homeomorphism on spectra
by Lemma \ref{editorial-source-lemma-p-ring-map}
and the same scalar comparison,
so it preserves and reflects
connectedness. Also finite Galois
extensions suffice in (2): every
finite separable $E/k$ embeds
in a finite Galois $N/k$, and
the injection $R\otimes_k E\to R\otimes_k N$
reflects nontrivial idempotents
and nonzeroness.

Finally, for every original $L/k$,
geometric connectedness of $R/k$
is equivalent to geometric
connectedness of $(R\otimes_k L)/L$.
The forward direction uses
the displayed scalar comparison
for any $F/L$. For the reverse
direction and any $E/k$, construct
a common receiving field for
$E$ and $L$ by the same
prime-quotient construction.
Geometric connectedness over
$L$ gives connectedness of
$R\otimes_k F$, and its
injected subalgebra $R\otimes_k E$
is nonzero and has no nontrivial
idempotents. This proves the
reverse implication with the
original scalar actions.
\end{proof}